\section*{Erd\H{o}s Problem 162: a logarithmic quasi-Ramsey threshold}
\subsection*{Statement and necessary corrections}
The page asks for $F(n,\alpha)\sim c_\alpha\log n$, but defines $F$ using the \emph{largest} threshold $k$ for which some colouring has more than $\alpha\binom{|X|}{2}$ edges of each colour in every $|X|\ge k$.
For fixed $\alpha<1/2$ and sufficiently large $n$, a balanced global colouring satisfies the condition at $k=n$, because
$\lfloor\binom n2/2\rfloor>\alpha\binom n2$.
Thus the literal largest threshold in $[1,n]$ is $n$; if larger thresholds are allowed there is no largest one. At $\alpha=1/2$ the strict requirement is impossible on any set of at least two vertices.

We explicitly study the intended \emph{smallest} threshold $Q(n,\alpha)$ for $0\le\alpha<1/2$, allowing $n+1$ as a vacuous value. The wording correction does not resolve the asymptotic question.

\subsection*{A correct clique lower bound}
Splitting a vertex's neighbours by incident colour gives
$R(a,b)\le R(a-1,b)+R(a,b-1)$, with
$R(1,b)=R(a,1)=1$. Induction consequently gives
\[
 R(r,r)\le\binom{2r-2}{r-1}\le4^{r-1}.
\]
Every colouring on $n$ vertices therefore has a monochromatic clique of size
$r=\lfloor\log_4 n\rfloor+1$. For large $n$, this is at least two and violates the strict requirement for every $\alpha\ge0$. Hence
\[
 Q(n,\alpha)\ge\lfloor\log_4 n\rfloor+2. \tag{1}
\]
The old local greedy proof allowed its majority colour to change at successive vertices and then incorrectly declared the chosen vertices monochromatic.

\subsection*{An explicit upper constant}
Put $I_\alpha=\log2+\alpha\log\alpha+(1-\alpha)\log(1-\alpha)$, with $0\log0=0$. A fixed $s$-set in an independent fair colouring is bad with probability at most
$2e^{-I_\alpha\binom s2}$. For $0<\alpha<1/2$, exponential Markov for the binomial red count, optimized at $e^{-z}=\alpha/(1-\alpha)$, proves this. For $\alpha=0$, the exact probability of a monochromatic clique is $2^{1-\binom s2}$.

For fixed $\eta>0$ take $s=\lceil(2/I_\alpha+\eta)\log n\rceil$. A union bound gives
\[
 \mathbb P(\text{some bad $s$-set})
 \le2\binom ns e^{-I_\alpha s(s-1)/2}<1
\]
for large $n$, since the coefficient of $(\log n)^2$ in its logarithm is negative. Summing the strict inequalities over all $s$-subsets of a larger set, and dividing by each edge's multiplicity, proves the condition at all larger sizes. Thus
\[
 \frac1{\log4}\le\liminf\frac{Q(n,\alpha)}{\log n}
 \le\limsup\frac{Q(n,\alpha)}{\log n}\le\frac2{I_\alpha}.
 \tag{2}
\]

\subsection*{The zero-density question contains a Ramsey limit problem}
For $n\ge2$ put $r(n)=\max\{r:R(r,r)\le n\}$. Every colouring has a monochromatic $r(n)$-clique, while some colouring has no $(r(n)+1)$-clique. Therefore
\[
 Q(n,0)=r(n)+1. \tag{3}
\]
A finite positive limit $Q(n,0)/\log n\to c$ is consequently equivalent to
$\log R(r,r)/r\to1/c$. If the latter limit exists, bracket $n$ between consecutive Ramsey thresholds in (3). Conversely, apply (3) at $n=R(r,r)$ and at $R(r,r)-1$, where the values of $r(n)$ are $r$ and $r-1$. The thresholds are strictly increasing: add one vertex, with arbitrary incident colours, to a colouring on $R(r,r)-1$ vertices avoiding monochromatic $r$-cliques. A monochromatic $(r+1)$-clique in the enlarged colouring would leave an old monochromatic $r$-clique, so $R(r+1,r+1)>R(r,r)$. The two evaluations prove the inverse limit assertion.

\subsection*{Assessment and attribution}
The random-colouring strategy of \texttt{gpt\_pro\_5.2/162.tex} is retained, with its false clique bound, endpoint treatment, and unchecked numerical claims replaced. Equations (1)--(3) are proved here. Existence of $c_\alpha$ for the intended threshold remains unresolved in this attempt. Source: \url{https://www.erdosproblems.com/162}. No novelty claim is made.
\subsection*{COMPLETION ESTIMATE}
\textbf{COMPLETION: 35\%}
